\documentclass[11pt]{article}
\usepackage[T1]{fontenc}
\usepackage[utf8]{inputenc}
\usepackage{lmodern,amsmath,amssymb,amsthm}
\usepackage[margin=1in]{geometry}
\usepackage[hidelinks]{hyperref}
\newtheorem{theorem}{Theorem}\newtheorem{lemma}[theorem]{Lemma}
\newcommand{\R}{\mathbb R}\newcommand{\Z}{\mathbb Z}
\newcommand{\C}{\mathbb C}\newcommand{\E}{\mathbb E}
\title{A normalized local converse for the GI assembly}
\author{Author proposal: /root/astra\_v14\_baseline\_theta}
\date{7 November 2026}
\begin{document}\maketitle

\noindent\textbf{Scope.} This author proof, awaiting independent review,
derives the precise correlation-to-Gowers-norm implication used by the
assembly. It imports the published compact cube constraint, Green--Tao,
\emph{Linear equations in primes}, Proposition 11.5, p.~1799, with its
definitions in 11.4 and proof in Appendix E. It uses the accepted local
increment linearization for polynomial tests. It does not import
Proposition 11.2's averaging-domain transition, assume an inverse
theorem, or claim to reprove the cube constraint.

\section{Available geometric assertion and needed conclusion}

Fix an integer $d\geq1$ and a fixed $d$-step nilmanifold $M=G/\Gamma$
with a smooth metric. The available cube constraint says there is a
compact subset
$\Sigma\subseteq M^{\{0,1\}^{d+1}\setminus\{0\}}$ and a continuous
map $P:\Sigma\to M$ such that, for every $g\in G$, $x\in M$,
$n\in\Z$ and $h\in\Z^{d+1}$, the nonzero vertices
$(g^{n+\omega\cdot h}x)_{\omega\ne0}$ lie in $\Sigma$ and
\[
 g^nx=P((g^{n+\omega\cdot h}x)_{\omega\ne0}).
\]
No bound on $g$ or the integers is imposed. In our application $G$ is
connected simply connected and at most $d$-step; these satisfy the
published nilmanifold hypothesis.

\begin{theorem}\label{thm:converse}
For fixed $d,M,L$ and every $\kappa>0$ there is $c>0$ such that
every $1$-bounded $v:[N]\to\C$ and $1$-bounded $L$-Lipschitz test
$F:M\to\C$ satisfying
\[
 |\E_{n\in[N]}v(n)F(g^nx)|\geq\kappa
\]
obey $\|v\|_{U^{d+1}[N]}\geq c$. The bound is uniform in $N,g,x,F$.
The same conclusion holds for a fixed finite family of rational
degree-at-most-$d$ polynomial models, using the exact same-step
increment linearization and taking a finite maximum of its bounds.
\end{theorem}

\section{A finite tensor approximation uniform in the test}

The algebra generated by coordinate functions that are Lipschitz on
$M$ separates points on the compact set $\Sigma$, contains constants
and is closed under conjugation. Complex Stone--Weierstrass gives,
for every fixed continuous $F$ and $\epsilon>0$, an approximation
\[
 F(P(z))=\sum_{a=1}^A c_a\prod_{\omega\ne0}H_{\omega,a}(z_\omega)
                 +R(z),\qquad |R(z)|\leq\epsilon,
\]
where the factors are $1$-bounded Lipschitz functions. Scaling each
factor into the unit disk is absorbed into the finite coefficients
$c_a$; no coefficient is discarded. By Arzela--Ascoli the class
$\|F\|_\infty\leq1$, $\operatorname{Lip}(F)\leq L$ is compact
in the uniform norm. A finite uniform net, approximated as above,
makes $A$, $\sum_a|c_a|$ and all factor bounds depend only on
$M,L,\epsilon$. Choose a net radius $\epsilon/2$ and approximate
its members to $\epsilon/2$ to retain the displayed total error.

For each orbit, put
$b_{\omega,a}(t)=H_{\omega,a}(g^tx)$. The cube constraint gives
\begin{equation}\label{eq:pointwise}
 F(g^nx)=\sum_a c_a\prod_{\omega\ne0}
 b_{\omega,a}(n+\omega\cdot h)+O(\epsilon)
\end{equation}
for every integer $n,h$, uniformly. The factors are $1$-bounded;
their dependence on the orbit will no longer matter in the analytic
estimate.

\section{Normalized smooth averaging and a uniform Fourier bound}

Choose once and for all a smooth $\sigma:\R\to[0,1]$ supported on
$[-1,2]$ and equal to one on $[0,1]$. Define
\[
 S_N=\sum_{h\in\Z}\sigma(h/N),\qquad
 \mu_N(h)=\sigma(h/N)/S_N.
\]
Then $N\leq S_N\leq3N+1\leq4N$, and $\mu_N$ is a probability
measure with finite support. Average \eqref{eq:pointwise} against
the product probability measure $\prod_j\mu_N(h_j)$ on all
integer $h$. The error remains at most $\epsilon$, with no missing
normalizing scalar. Thus
\[
 F(g^nx)=\sum_a c_a B_a(n)+O(\epsilon),\quad
 B_a(n)=S_N^{-(d+1)}\sum_{h\in\Z^{d+1}}
   \prod_j\sigma(h_j/N)\prod_{\omega\ne0}
                  b_{\omega,a}(n+\omega\cdot h).
\]

Take $Q=C_dN$ with a fixed integer $C_d>10(d+2)$, and identify
representatives of $\Z/Q\Z$ with a centered interval. The sampled
cutoff $w_N(h)=\sigma(h/N)$ has support strictly inside that interval.
Write its discrete Fourier expansion as
$w_N(h)=\sum_r\widehat w_N(r)e(rh/Q)$ with normalized coefficients.
There is a bound
\begin{equation}\label{eq:l1}
 \sum_{r\in\Z/Q\Z}|\widehat w_N(r)|\leq C'_d
\end{equation}
independent of $N$. To prove it, the zero coefficient and every
coefficient are at most $S_N/Q=O_d(1)$. For $r\ne0$, summation by
parts twice gives
\[
 |1-e(r/Q)|^2 |\widehat w_N(r)|
 \leq Q^{-1}\sum_h|\Delta^2w_N(h)|\leq C/(NQ).
\]
The final bound follows from the second derivative of the fixed smooth
cutoff, with only $O(N)$ possibly nonzero terms. For the distance
$r_*\in[1,Q/2]$ to zero modulo $Q$,
$|1-e(r/Q)|\geq4r_*/Q$. Thus
$|\widehat w_N(r)|\leq C_d/r_*^2$ and summation proves
\eqref{eq:l1}. Smooth zero extension of $\sigma$ ensures that the
same finite-difference bound holds at its support boundaries.

\section{The exact Gowers estimate}

Zero-extend $v$ to $\Z/Q\Z$. Extend each bounded integer function
$b_{\omega,a}$ to the cyclic group so that it agrees on all integer
arguments that occur when $n\in[N]$ and $h_j\in[-N,2N]$.
Those arguments lie between $-(d+1)N$ and $(2d+3)N$; the choice
of $C_d$ places the whole interval strictly inside the centered
representative range. There is therefore no conflicting cyclic value.
The normalized correlation with $B_a$ is exactly
\[
 \frac{Q}{N}\left(\frac{Q}{S_N}\right)^{d+1}
 \E_{n,h\in\Z/Q\Z}
 v(n)\prod_j w_N(h_j)
          \prod_{\omega\ne0}b_{\omega,a}(n+\omega\cdot h).
\]
The prefactor is bounded above by a constant depending only on $d$.
This identity comes from an all-integer smooth average; no hard box
average is replaced by a larger summation domain.

Expand the product of cutoffs by Fourier series. The sum of absolute
values of its coefficients is at most $(C'_d)^{d+1}$. In each term,
the factor $e(\sum_jr_jh_j/Q)$ can be absorbed into the nonzero
vertex functions. For each $j$, choose an index $k\ne j$, which
exists since $d+1\geq2$, and use
$h_j=(n+h_j+h_k)-(n+h_k)$. Multiply the functions at vertices
$e_j+e_k$ and $e_k$ by the corresponding opposite linear phases.
Their absolute bounds remain one. The Gowers--Cauchy--Schwarz
inequality for the full $(d+1)$-cube gives the absolute bound
$\|v\|_{U^{d+1}(\Z/Q\Z)}$ for every such term. Complex
conjugations are absorbed into the vertex functions, leaving the
norm and absolute bounds unchanged. Hence
\[
 |\E_{[N]}v B_a|\leq K_d\|v\|_{U^{d+1}(\Z/Q\Z)}
 \leq K_d\|v\|_{U^{d+1}[N]}.
\]
The last inequality uses the prescribed normalized interval norm:
the cyclic norm is the interval norm multiplied by the cyclic norm
of the interval indicator, which is at most one. If necessary enlarge
$C_d$ also to exceed $2^{d+1}$ so this single ambient group realizes
that interval normalization.

Combining the finite tensor approximation with this estimate yields
\[
 |\E_{[N]}v(n)F(g^nx)|
 \leq K_d\Bigl(\sum_a|c_a|\Bigr)
                     \|v\|_{U^{d+1}[N]}+\epsilon.
\]
Choose $\epsilon<\kappa/2$. This proves
Theorem~\ref{thm:converse} with a positive bound depending only on
$d,M,L,\kappa$. A fixed bound other than one for $v$ or $F$ is
handled by dividing by that bound first. This proof is complex-valued
throughout; it does not need a real-valued dual-norm convention.

\section{Polynomial input and source comparison}

For the actual assembly input, a scalar polynomial nilsequence has
degree $d=s-2\geq1$ on fixed rational geometry. The accepted
increment linearization001 represents it exactly on a fixed group
of step at most $d$, with a fixed Lipschitz multiplier and unchanged
norm bound. Applying the theorem there and taking finite maxima
over a finite family gives precisely the positive normalized
$U^{s-1}[N]$ norm required by the last induction step. This is a
converse implication and has no inverse-theorem premise.

The literal published proof of 11.2 on p.~1801 defines an average
over $h\in[N]^{d+1}$ with $\sigma(h_j/N)$, while $\sigma=1$ on
$[0,1]$. Its later passage from (11.6) to (11.7) widens the shift
domain to a cyclic group with the same cutoff supported on $[-1,2]$.
That changes the set of summands. The local proof above instead
starts with the normalized all-integer smooth average and tracks
its exact cyclic prefactor. Whether an earlier source correction
already supplied this interpretation has not been established.
The observation is a bounded proof-transition issue, not a
counterexample to Proposition 11.2 or the converse theorem.

The related-use scope is the normalization of the 11.2 average,
its dual-norm estimate, Corollary 11.6, and the present GI converse
application. Published 11.5 remains an explicit unchanged supplier
whose source-level application must be independently accepted.
No claim of whole-source acceptance of Appendix E or of the full
GI assembly follows from this local analytic proof.

\end{document}
